% ------------------------------------------------------------
% Page 91
% ------------------------------------------------------------

Heureusement la famille Hubac a conservé un certain nombre d’archives à partir de 1866, en
particulier les actes notariés des ventes successives, Connillergues passant des mains d’une
famille \textbf{Angliviel}, probablement propriétaire-exploitant, à un certain \textbf{Benjamin Salles}, rentier
et propriétaire habitant Valleraugue, le 3 Novembre 1866. Le prix avait été fixé à 20.000
francs

C’est en 1897 que sera construit le pont actuel qui remplace une passerelle située à une
vingtaine de mètres en aval et dont a subsisté très longtemps, un pieu de châtaignier, au milieu
de la rivière.

\textbf{En 1919} un petit fils de Madame Salle, qui avait perdu son mari, vendit Connillergues à \textbf{Jules
Laget} pour 40.000 francs. Domaine affermé depuis 1911 jusqu’à 1920 à Albert Gaches de
Rousses, puis des ouvriers agricoles font « tourner » l’exploitation. (l’hiver ils descendent
dans les vignes du Languedoc) jusqu’à l’arrivée de Jules et Irène Hubac, fermiers, en 1947.
Jules.Laget revend rapidement, \textbf{en 1923 à Louis Hubac}, propriétaire, né à Gatuzieres dit « le
Huguenâot » car il avait été fermier à Sérigas où il était le seul protestant, Il est probable qu’il
descend de ces Hubac dont nous trouvons la trace à Gatuzieres en 1687 sur les registres des
Nouveau Convertis au catholicisme:

\textit{« Hubac et Marguerite Fages sa femme. 60 ans, ils ont 3 enfants. La fille aînée fait assez bien.
Marguerite : mal. Pierre très mal ; scandaleux, tient des discours impies et d’un scélérat,
tenant une posture indigne à l’église ; étant à craindre qu’il n’aille aux assemblées s’il s’en
faisait ».}

Selon Elie Chanson, ancien fermier (1958-1979) la vente se serait déroulée rapidement,
oralement, dans les rues de Meyrueis, pour la somme de 30.000 francs, aussitôt attestée par un
compromis signé devant notaire. Le lendemain le vendeur se serait rétracté, mais Louis Hubac
aurait fait valoir le compromis, signé la veille, pour contraindre la vente.

\textbf{1947} Jules et Irène Hubac, née Maurin, mariés au mois d’Août, habitent Connillergues et y
travaillent probablement sans bail de fermage lequel ne sera établi qu’en 1950 et pour 9 ans
entre

\underline{Propriétaires} : Louis Paul Hubac, pasteur à Roquecourbe

\hspace{3.5cm} Adeline Hubac, veuve de Eugène Maurin sa sœur

\hspace{3.5cm} Marceau Hubac, son frère.

\underline{Fermiers} : Jules et Irène Hubac (à noter qu’Irène était la nièce d’Eugène Maurin).

\vspace{0.2cm}

Fermage annuel payable par semestre à la valeur de :

\hspace{3.5cm} 1807 litres de lait de vache

\hspace{3.5cm} 228 kg de viande de veau

A cette époque Connillergues accueillait l’été, en estive, un « transhumant » de 1200 à 1300
brebis. Jules et Irène Hubac sont restés fermiers de 1947 à 1954.
André Hubac, fils de Louis, un des propriétaires, « élevait » des abeilles avec l’aide de
Paparel, coiffeur à la Villa des Roses et apiculteur. Les ruches étaient placées près du
pigeonnier et Jules Hubac s’y intéressait de près, ne craignant pas les piqûres et n’hésitant pas
à rechercher et ramasser les essaims.
